% ------------------------------------------------------------------
%  Computing in Cardiology 2026 -- draft manuscript
%  Build: build.bat  (xelatex -> bibtex -> xelatex x2)
%  NOTE: the official CinC class (cinc.cls) is not installed on this
%  machine; the layout below imitates it with the article class.
%  Replace the preamble by the official template before submission.
% ------------------------------------------------------------------
\documentclass[10pt,twocolumn,a4paper]{article}
\usepackage[utf8]{inputenc}
\usepackage[T1]{fontenc}
\usepackage{newtxtext,newtxmath}
\usepackage[a4paper,top=2.3cm,bottom=2.3cm,left=1.9cm,right=1.9cm,columnsep=0.6cm]{geometry}
\usepackage{graphicx}
\usepackage{booktabs}
\usepackage{array}
\usepackage{amsmath}
\usepackage{caption}
\usepackage{titlesec}
\usepackage{microtype}
\usepackage[hidelinks]{hyperref}

\pagestyle{empty}
\setlength{\parindent}{1.2em}
\setlength{\parskip}{0pt}
\captionsetup{font=small,labelfont=bf,skip=3pt}
\titleformat{\section}{\normalfont\large\bfseries}{\thesection.}{0.5em}{}
\titleformat{\subsection}{\normalfont\normalsize\bfseries}{\thesubsection.}{0.5em}{}
\titlespacing*{\section}{0pt}{6pt plus 2pt minus 2pt}{2pt}
\titlespacing*{\subsection}{0pt}{4pt plus 2pt minus 2pt}{1pt}
\setlength{\tabcolsep}{3.0pt}
\newcommand{\fitwidth}[1]{\resizebox{\ifdim\width>\columnwidth\columnwidth\else\width\fi}{!}{#1}}
\renewcommand{\arraystretch}{1.05}
\newcommand{\ms}{\,ms}
\newcommand{\Hz}{\,Hz}
% tighter reference list (the official CinC template sets references solid)
\let\oldthebibliography\thebibliography
\renewcommand{\thebibliography}[1]{\oldthebibliography{#1}%
  \setlength{\itemsep}{0pt}\setlength{\parsep}{0pt}\setlength{\parskip}{0pt}}

\title{\vspace{-1.2em}\Large\bfseries Single-Channel Fetal QRS Detection:\\ Where the Performance Actually Comes From}
\author{Ngo Binh Minh$^{1}$, [Co-author]$^{1}$, [Supervisor]$^{1}$\\[3pt]
\normalsize $^{1}$Faculty of Information Technology, Ton Duc Thang University, Ho Chi Minh City, Vietnam}
\date{}

\begin{document}
\twocolumn[
\begin{@twocolumnfalse}
\maketitle
\vspace{-2.0em}
\begin{center}
\begin{minipage}{0.92\textwidth}
\section*{\centering\normalsize Abstract}
\small\itshape
Deep detectors for fetal QRS in single-channel abdominal ECG are reported as a whole, so it is
unclear which design choices produce the accuracy. We hold the protocol fixed (leave-one-record-out,
$\pm$50\ms, threshold from an inner validation record, label-blind lead selection), vary one factor
at a time, and estimate every effect per woman. A dilated temporal convolutional network with
113\,481 parameters reaches a macro F1 of 99.21\% (95\% CI 97.9--100.0) on ADFECGDB and 93.30\%
(84.9--99.4) zero-shot on ten pregnancy recordings of the Silesian database. A three-tier
decomposition attributes the gain to the front-end (the 10--60\Hz\ band of Xu et al., $+$11.00 F1,
on a 300\ms\ window learner), to context and per-sample output ($+$4.49, CI 1.6--7.4) and to the
architecture family ($+$0.41, not resolvable at $n=5$). Repeated on the TCN itself over 22 women the same band
change is worth only $+$2.44 (CI {$[-0.04;6.29]$}), three women carrying all of it, so tier 1 does
not transfer; that run is preliminary. Re-run by us under GNU Octave on all 22
women with the same matcher, the Power-MF pipeline gains 12.12 F1 points from its four channels over
its own one-channel restriction (86.71 to 98.83); the network recovers 10.85 of those points
(89.5\%) from one lead, so it is not separable from four-channel Power-MF ($-$1.27 F1, CI
{$[-3.08;+0.27]$}, $p=0.16$, better in 18 of 22 women, median $+$0.23) and is ahead of one-channel
Power-MF ($+$10.85, 22 of 22). On all 75 CinC 2013 set-a records, training on 22 women instead of
five raises the label-blind score from 71.21 to 79.40 ($p=3\cdot10^{-10}$); the ten-record sample of
earlier drafts was biased by $-$12 points under the blind rule and $+$19 under a post-hoc lead, and
both figures are withdrawn. Short-term variability inherits detection error one-sidedly (bias
$+$20.50\ms\ on CinC 2013) and is not a standalone variability monitor.
\end{minipage}
\end{center}
\vspace{1.2em}
\end{@twocolumnfalse}
]

\section{Introduction}
Fetal heart rate from abdominal electrodes needs a reliable fetal QRS (fQRS) detector; a
single-channel one keeps the hardware small. The PhysioNet/CinC Challenge 2013 fixed the metric
still in use, F1 with a $\pm$50\ms\ window \cite{clifford2014}; Behar et al.\ benchmarked template
subtraction, Kalman and blind-source methods under it, grid-searching the band-pass filter
\cite{behar2014}; Andreotti et al.\ showed on synthetic signals that pre-processing changes the
ranking of extraction algorithms \cite{andreotti2016}. Deep single-channel methods since
\cite{zhong2018} change many things at once: dual-path source separation \cite{shokouhmand2023}, a
matched-filter pipeline with power-spectral-density (PSD) lead selection over four channels
\cite{jaeger2024}, and guided empirical mode decomposition with a 10--60\Hz\ band \cite{xu2026}.

Several ingredients of our detector are established elsewhere and we claim none as new: per-sample
R-peak output instead of window classification \cite{zahid2022}, dilation rather than kernel size as
the useful knob in fECG \cite{fotiadou2021}, and the band itself \cite{xu2026}. Each of those papers
changes band, sampling rate, context, architecture and sometimes the lead rule at once, and the lead
is often chosen with the test labels; what is missing is how much each choice is worth once the
protocol is held fixed. The contributions are:
\begin{itemize}\setlength{\itemsep}{0pt}
\item a leakage-free single-channel protocol with a three-tier decomposition of the gain into
  front-end, temporal context and architecture family, analysed per woman;
\item a like-for-like comparison against a multi-channel published pipeline re-run by us and scored
  with the same matcher, rather than quoted from its paper;
\item cross-dataset evidence on 22 independent women and all 75 CinC 2013 set-a records that
  generalisation breaks across recorders rather than gestation, and the label-blind lead rule with
  it.
\end{itemize}

\section{Methods}
\subsection{Data}
\textbf{ADFECGDB} \cite{goldberger2000}: five labour recordings (r01, r04, r07, r08, r10), four
abdominal leads at 1\,kHz, 5\,min each, with fQRS annotations from a scalp electrode (3\,191 beats).
\textbf{Silesian database} \cite{matonia2020}: ten pregnancy recordings (B1, 32--42 weeks,
199.6\,min total) and twelve labour recordings (B2, 5\,min), four abdominal leads at 500\,Hz, same
hospital and recorder as ADFECGDB. B1 annotations are indirect (author pipeline plus expert
correction) and in 5/10 records sit 8--12\ms\ after the abdominal R peak; B2 annotations come from a
scalp electrode. The overlap with ADFECGDB was noted by Baldazzi and Pani \cite{baldazzi2023}; we
cross-correlated all $4\times4$ lead pairs at 250\,Hz and found five B2 records that \emph{are}
ADFECGDB records (B2\_01\,=\,r01,
B2\_02\,=\,r10, B2\_07\,=\,r04, B2\_10\,=\,r07, B2\_11\,=\,r08; correlation 0.988--0.994). The
independent set is therefore 22 women
(5 ADFECGDB, 10 B1, 7 new B2); duplicates are scored only with the fold model that never saw
them. \textbf{CinC 2013 set-a}
\cite{clifford2014}: all 75 records, four channels at 1\,kHz, 1\,min each, a different recorder.

\subsection{Detector}
One abdominal lead is band-passed 10--60\Hz\ (4th-order Butterworth, zero-phase) with a 50\Hz\ notch
and resampled to 250\,Hz; the band is taken from \cite{xu2026} and Sec.~3.4 measures what it is
worth against seven other published bands. Maternal QRS are located on an 8--25\Hz\ envelope and
removed by subtracting a median beat template with per-beat least-squares scaling. The residual
and the filtered signal, each robust-scaled, form a two-channel 4\,s input. The network is a dilated
residual TCN: a stem and five residual blocks (dilations 1, 2, 4, 8, 16; 40 channels; kernel 7;
BatchNorm; GELU) with a $1\times1$ output convolution, 113\,481 parameters, receptive field
1\,516\ms. It emits one logit per sample against a Gaussian heat-map target ($\sigma=3$ samples) at
each annotated fQRS \cite{zahid2022}, trained with weighted binary cross-entropy and AdamW
($3\times10^{-3}$, one-cycle), six epochs, batch 32, 1\,s stride. At inference, overlapping windows
are averaged and peaks above a threshold with a 250\ms\ refractory period are the fQRS (0.48\,MB,
4.35\ms\ per window on a desktop CPU).

\subsection{Protocol and statistics}
Leave-one-record-out (LORO) on ADFECGDB with five folds: one test, one inner validation and three
training records; the threshold is chosen on the validation record from a grid 0.20--0.75 and
applied unchanged. Detections are matched one-to-one to
annotations by greedy nearest matching within $\pm$50\ms. Because a device sees one lead, we also
report a \emph{label-blind} lead: following Power-MF \cite{jaeger2024}, the lead whose
residual envelope has the largest Welch PSD peak in 1.8--3.0\Hz, a rule that never touches
annotations. Silesian and CinC records are scored zero-shot. A second model, used only where stated,
is trained on all 22 independent women with 11-fold grouped cross-validation.

The unit of analysis is the woman, not the (record$\times$lead) pair: four leads of one recording
are repeated measures in one cluster, and treating them as 20 observations inflates every test. We
average leads within a woman and report mean differences with 95\% cluster-bootstrap percentile
intervals over subjects (10\,000 resamples, seed 0); where these disagree with the conservative
paired-$t$ interval at $n\le7$ we give both and trust the latter. ``No difference'' claims use two
one-sided tests (TOST) at a pre-declared 1.0-point margin. With $n=5$ the smallest two-sided $p$ an
exact signed-rank test can give is $0.0625$, so no ADFECGDB comparison can reach $p<0.05$ at subject
level; we quote intervals there, not $p$-values.

\subsection{Experiments}
\textbf{Baselines}, re-implemented and scored with the same matcher: template subtraction (TS) and
TS with a two-component principal-component template (TS-PCA) \cite{behar2014}, each with an
adaptive Pan--Tompkins detector \cite{pan1985}, and \emph{Prominence}, our front-end followed by
\texttt{scipy} peak picking; hyper-parameters were tuned on r01 only. \textbf{Power-MF}
\cite{jaeger2024} was not quoted but executed: the authors' MATLAB pipeline was run under GNU Octave
11.3 and scored with \emph{our} matcher at $\pm$50\ms, at four leads with both ICA stages and again
restricted to the single PSD-selected lead, so all three columns share a scoring rule.

\textbf{Three-tier decomposition.} Tier 1 (front-end): eight literature band-pass settings, LORO,
gradient-boosted trees on 300\ms\ windows of two leads. Tier 2: a 300\ms\ window classifier versus
the 4\,s per-sample TCN. Tier 3 (family at fixed context): gradient boosting versus a dilated CNN on
the same windows. A receptive-field sweep (172--3\,052\ms) isolates context within one family, and
an architecture benchmark compares eight families on 300\ms\ windows.

\textbf{Reject gate.} A 4\,s segment is \emph{bad} if the detector's F1 on it is below 80. A
classifier trained on 1\,500 ADFECGDB segments and tested on 600 CinC 2013 segments uses 16
topological \cite{adams2017}, 12 classical \cite{andreotti2017}, or all features; intervals resample
records, not segments.

\section{Results}
\subsection{Detection across data configurations}
Table~\ref{tab:data} gives the same model under the same protocol on five configurations. On
ADFECGDB the label-blind PSD rule reaches 99.21\%, identical to the oracle lead in every record, and
99.18\% on the DPSS benchmark leads \cite{shokouhmand2023}; across four leads the
LORO macro F1 is 97.45\%, 97.16\% over three seeds. Zero-shot on
Silesian pregnancy records it keeps 93.30\%, eight of ten records above 90 (oracle lead 97.04\%); the
two weakest are those where the PSD rule picks the worst lead. A second seed of the 22-subject model moves
per-subject F1 by 0.28 points on average (largest 2.81), though per-fold thresholds move more.

\begin{table}[t]
\centering\footnotesize
\caption{Macro F1 (\%) of one model trained on five ADFECGDB women, $\pm$50\ms; CI is a 95\% cluster
bootstrap over records, ``$\ge$90''/``$<$50'' count PSD-lead F1 in that range.}
\label{tab:data}
\fitwidth{\begin{tabular}{@{}lcccccc@{}}
\toprule
Configuration & $n$ & PSD lead & 95\% CI & 4-lead mean & $\ge$90 & $<$50\\
\midrule
ADFECGDB, LORO & 5 & 99.21 & [97.9; 100.0] & 97.45 & 5 & 0\\
Silesia B2 labour, all & 12 & 97.17 & [92.7; 99.8] & 95.26 & 11 & 0\\
\quad 7 new women, zero-shot & 7 & 95.87 & [88.4; 99.8] & 93.73 & 6 & 0\\
Silesia B1 pregnancy, zero-shot & 10 & 93.30 & [84.9; 99.4] & 91.31 & 8 & 0\\
CinC 2013 set-a, zero-shot & 75 & 71.21 & [63.6; 78.4] & 64.78 & 42 & 22\\
\bottomrule
\end{tabular}}
\end{table}

\subsection{How much of multi-channel separation one channel recovers}
Competing numbers in the previous draft were quoted, not measured, so we ran Power-MF
\cite{jaeger2024} under Octave and scored it with our matcher. That draft reported
a comparison over 16 of 22 women because the port returned nothing, or half the beat rate, on six
long B1 recordings. \emph{Those figures are withdrawn.} The cause was a defect in our port, not a
property of the method: Octave's \texttt{signal} \texttt{findpeaks} enforces
\texttt{MinPeakDistance} with a pairwise distance matrix, $O(k^{2})$ in memory, exceeding the index
limit on B1 records (2\,395\,600 samples after the pipeline's $\times4$ interpolation), whereas
MATLAB --- the authors' platform --- resolves it greedily in $O(k\log k)$. We
re-implemented the MATLAB semantics with a doubly linked list ($O(k)$, agreeing on 48/48 test cases).
The earlier explanation, that the 340\ms\ minimum peak distance sat under the fetal RR interval, was
tested against the reference annotations and rejected: 0.00\% of the RR intervals of B1\_01 are
shorter than 340\ms. Externally, the repaired port scores 99.40\,$\pm$\,0.51 on the ten B1 records
against 99.46 published by the authors, a gap of 0.06. All 22 women are scored, with no exclusion
(Table~\ref{tab:pmf}).

This makes the question explicit: how much of what \emph{multi-channel source separation} buys can
a 113\,k single-channel network recover? Restricting Power-MF to one lead costs it 12.12 F1
points (98.83 to 86.71, CI $[-18.27;-6.90]$, worse in 21 of 22 women), which is the price of the
four-channel ICA stages on this material. Our single-channel network recovers 10.85 of those points
($+$10.85, CI $[+6.80;+15.40]$, $p=5\cdot10^{-7}$, better in 22 of 22 women), that is 89.5\% of the
benefit, from one electrode pair. Against the full four-channel pipeline it is not separable
($-$1.27, CI $[-3.08;+0.27]$, $p=0.16$), and on ADFECGDB the interval excludes zero in its favour
($+$0.39, CI $[+0.22;+0.60]$, 5/5 women).

The mean difference against four channels is negative, but it is not a uniform deficit: the median
difference over the 22 women is $+$0.23 and the network is better in 18 of them. The mean is pulled
below zero by three recordings --- B1\_07 ($-$12.88), B1\_06 ($-$10.51) and B2\_03 ($-$9.77) --- and
mildly by B1\_10 ($-$1.86). The defensible statement is that one channel recovers most of the
multi-channel benefit, not that it beats four channels.

\begin{table}[t]
\centering\footnotesize
\caption{Power-MF \cite{jaeger2024} \emph{re-run} by us, at four channels and restricted to one
lead, against our 22-subject single-channel model: same recordings, same $\pm$50\ms\ matcher.
$\Delta_{4}$ and $\Delta_{1}$ are ours $-$ Power-MF per woman with 95\% cluster-bootstrap intervals
(10\,000 resamples of subjects). Power-MF was not run at four channels on CinC 2013; the CinC
$\Delta_{1}$ is a mean difference, without an interval.}
\label{tab:pmf}
\fitwidth{\begin{tabular}{@{}lcccccc@{}}
\toprule
 & & \multicolumn{3}{c}{mean F1 (\%)} & \multicolumn{2}{c}{ours $-$ Power-MF}\\
\cmidrule(lr){3-5}\cmidrule(lr){6-7}
Subject set & $n$ & PMF-4 & PMF-1 & ours & $\Delta_{4}$ [95\% CI] & $\Delta_{1}$ [95\% CI]\\
\midrule
ADFECGDB & 5 & 99.01 & 92.37 & 99.40 & $+0.39$ [$+0.22$; $+0.60$] & $+7.02$ [$+3.84$; $+9.34$]\\
Silesia B2, new & 7 & 97.90 & 87.28 & 96.82 & $-1.08$ [$-4.04$; $+0.54$] & $+9.54$ [$+3.09$; $+17.9$]\\
Silesia B1 & 10 & 99.40 & 83.48 & 97.15 & $-2.25$ [$-5.55$; $+0.32$] & $+13.68$ [$+6.91$; $+21.3$]\\
\textbf{All 22} & \textbf{22} & \textbf{98.83} & \textbf{86.71} & \textbf{97.56} & $\mathbf{-1.27}$ [$-3.08$; $+0.27$] & $\mathbf{+10.85}$ [$+6.80$; $+15.4$]\\
\midrule
CinC 2013 set-a & 75 & --- & 62.82 & 79.40 & --- & $+16.58$\\
\bottomrule
\end{tabular}}
\end{table}

\subsection{All 75 CinC 2013 records, and what ten records hid}
Earlier drafts scored ten of the 75 set-a records. Table~\ref{tab:cinc} reports the full set for the
five-woman (m5) and 22-woman (m22) models, on all 75 records and on the 68 left after removing a33,
a38, a47, a52, a54, a71 and a74, whose reference annotations are known to be unreliable
\cite{behar2014} (the list was declared before the run). Under the label-blind PSD rule, training on
22 women instead of five raises the mean from 71.21 to 79.40 ($+$8.18, CI $[5.62;11.09]$,
$p=3.2\cdot10^{-10}$, better in 53 records and worse in 5, Cliff's $\delta=0.23$); records at or
above 90 go from 42 to 48 and those below 50 from 22 to 16, and the 68-record variant gives $+$8.62
($p=3.7\cdot10^{-9}$).

The ten-record sample was biased in \emph{both} directions, which is why the complete set matters.
Under the blind rule it was pessimistic by about 12 points (m5 59.15 vs 71.21; m22 69.31 vs 79.40);
under a fixed lead 0 --- a post-hoc choice made after the PSD rule had been seen to fail --- it was
optimistic by 19--21 points (m5 77.34 vs 58.72; m22 90.34 vs 69.33). On the full set that post-hoc
lead is in fact the \emph{worst} of the four rules, ten points below the blind rule it was meant to
rescue; both ten-record figures are withdrawn. The blind-to-oracle gap also shrinks, from
22.33 on the ten records to 7.47 on all 75, but stays four times the in-domain gap (0.02 on ADFECGDB, 0.66 on
B2, 1.90 on B1).

\begin{table}[t]
\centering\footnotesize
\caption{All 75 CinC 2013 set-a records, zero-shot: mean F1 (\%) after training on five women (m5)
and on 22 (m22). $\Delta$ is per record, 95\% bootstrap; $p$ exact signed-rank; ``$\ge$90''/``$<$50''
count m22 records. Only the top row is label-blind.}
\label{tab:cinc}
\fitwidth{\begin{tabular}{@{}lccccccc@{}}
\toprule
Lead rule & $n$ & m5 & m22 & $\Delta$ [95\% CI] & $p$ & $\ge$90 & $<$50\\
\midrule
PSD, label-blind & 75 & 71.21 & \textbf{79.40} & $+8.18$ [5.6; 11.1] & $3\cdot10^{-10}$ & 48 & 16\\
4-lead mean & 75 & 64.78 & 74.09 & $+9.31$ [7.0; 11.8] & $4\cdot10^{-11}$ & 30 & 17\\
Fixed lead 0 (post-hoc) & 75 & 58.72 & 69.33 & $+10.61$ [7.5; 14.0] & $1\cdot10^{-8}$ & 36 & 27\\
Oracle lead (labels) & 75 & 78.21 & 86.87 & $+8.66$ [5.9; 11.7] & $1\cdot10^{-9}$ & 53 & 8\\
\midrule
PSD, 68 records & 68 & 72.08 & \textbf{80.70} & $+8.62$ [5.7; 11.8] & $4\cdot10^{-9}$ & 46 & 14\\
\bottomrule
\end{tabular}}
\end{table}

\subsection{Where the performance comes from}
Three tiers decompose the gain on ADFECGDB (LORO, $\pm$50\ms; tier 1 on two leads with a
gradient-boosted learner on 300\ms\ windows, tiers 2--3 on four leads). \emph{Tier 1, front-end}:
changing only the band from 1--45\Hz\ to 10--60\Hz\ moves F1 from 80.06 to 91.06, $+$11.00 points
(not stored per record, so no interval); the spread over eight literature bands is 18.89 points
(72.17 for 0.5--100\Hz\ to 91.06 for 10--60\Hz). Re-run on the TCN itself over the 22 women
(three folds, one seed, four epochs) the same change gives 95.63 to 98.07, $+$2.44, cluster
bootstrap $[-0.04;6.29]$ containing zero: at $n=22$ the band effect is \emph{not} established on the
model we use. The median woman moves 0.00; B2\_03 ($+$37.8), B1\_06 ($+$8.8) and B1\_07 ($+$8.2)
carry all of it, the other 19 average $-$0.07, and averaging the four leads instead of selecting one
leaves $-$0.07 $[-0.51;0.39]$ --- the gain acts through the lead rule, not the band. \emph{Tier 2, context and output}: a 300\ms\ window
classifier to the 4\,s per-sample TCN, 92.67 to 97.16, $+$4.49 per woman [1.58; 7.43] (the
conservative paired-$t$ interval $[-0.29;9.27]$ still contains zero at $n=5$); this step bundles
receptive field, output parameterisation and parameter count (33\,k to 113\,k). \emph{Tier 3,
family}: gradient boosting to a dilated CNN at the same context, 92.49 to 92.90, $+$0.41, inside the
1.0-point equivalence margin and not resolvable either way. Within the TCN family, LORO F1 rises
from 92.67 at a 172\ms\ receptive field to 97.16 at 1\,516\ms, and the validation-record sweep
saturates near 1.5\,s.

Against re-implemented baselines on ADFECGDB (four-lead F1: TS 78.96, TS-PCA 91.05, Prominence
86.39, ours 97.45) the network gains 11.06 points over peak picking on the residual it
receives (CI $[7.80;14.51]$, 5/5) and halves the spread (SD 4.44 vs 10.17). The 18.49-point margin
over TS is unstable (paired-$t$ $[-1.81;38.80]$, one record drives it), and TS-PCA on the PSD lead
(96.74) is 2.47 points below.

In the architecture benchmark six families lie within 1.53 points (92.43 to 90.90); TOST at the
1.0-point margin declares 0 of 7 equivalent and 6 of 7 inconclusive, so the absent difference
measures the sample size, not the architectures. Only two effects clear it: the plain CNN, same
parameters without dilation, is worse by 2.95 (90\% CI $[-4.55;-1.47]$), the large CNN 0.73 above
($[0.18;1.22]$). All eight, at 300\ms, stay below the 4\,s per-sample model.

\subsection{Reject gate, topology, clinical read-out}
The 12 classical features predict bad segments with AUROC 0.929 (CI $[0.830;0.982]$) and lift
retained-segment F1 from 61.96 at full coverage to 69.40 at 80\% (random 62.03). That pooled figure
is largely \emph{between-record}: over the five records containing both classes the
\emph{within-record} AUROC is only 0.721 ($[0.517;0.898]$), and that is what a monitoring session
sees. Topology contributed nothing: 16 topological features reach 0.566 and adding them
\emph{lowers} AUROC to 0.905.

\textbf{Clinical read-out.} Short-term variability (STV, 3.75\,s epochs) computed from the detected
beats tracks detection quality almost exactly: bias $+$0.33\ms\ on ADFECGDB and $+$20.50\ms\ on CinC
2013; stratified by F1, mean absolute STV error is 0.13\ms\ at F1\,$\ge$\,99.5 and 25.78\ms\ at
F1\,$<$\,90. The error is one-sided (25 of 32 records overestimate): every missed or spurious beat
perturbs an RR interval and the perturbations only add. On gated records the limits of agreement are
$[-0.39;+0.61]$\ms, but no in-domain subject has a reference STV near the range where the
measurement would change management, so the detector is usable for trending on gated records and
\emph{not} as a standalone STV monitor.

\section{Discussion}
\textbf{The front-end dominates the 300\ms\ learner, not yet the network.} Of the 15.9 points
gained over the three tiers, 11.0 come from the pass-band alone on that learner: the quantitative
form, on real recordings, of what \cite{andreotti2016} showed on synthetic ones. On the network the
same change is worth 2.44, interval containing zero, and nothing at all when the four leads are
averaged instead of selected. Architecture comparisons run under different
bands are still not comparisons of architectures. The receptive field is the architectural variable we can resolve: dilation ($+$2.95
at equal parameters) and a longer field ($+$4.5 from 172 to 1\,516\ms) clear the 1.0-point margin
\cite{fotiadou2021}; other family differences do not.

\textbf{One channel recovers most of what four channels buy.} Multi-channel separation is worth a
directly measured 12.12 F1 points to Power-MF itself, and a 113\,k single-channel network recovers
10.85 of them, 89.5\%. The residual gap against four channels is concentrated, not diffuse: three of
22 recordings account for it and the median woman is $+$0.23. Reducing the electrode count
costs little on average here, but the failures it creates are severe individually --- what a device
must gate against.

\textbf{Generalisation breaks across recorders, and the lead rule with it.} The model loses
about four points on pregnancy records from the same recorder (97 to 93) and about 26 on a different
one (97 to 71 under the blind rule). Part of that is the selection rule: the PSD rule is within 1.90
points of the oracle in domain and 7.47 below it out of domain. Training on 22 women recovers 8.2
points without closing the recorder gap. A learned gate should replace the spectral heuristic; ours
recovers 7.4 points at 80\% coverage between records, though within-record discrimination (0.721)
is what a bedside sees.

\textbf{Limitations.} The binding limitation is $n$: five women on ADFECGDB and 22 in all,
one hospital, one recorder, every training woman in labour at 38--41 weeks. No ADFECGDB comparison
can reach $p<0.05$ at subject level, every interval is wide, and the architecture benchmark is
underpowered by construction. \emph{Against four-channel Power-MF our mean
is 1.27 points lower.} The interval $[-3.08;+0.27]$ contains zero, so neither system can be declared
better, but the sign of the point estimate is against us and we do not claim otherwise; the deficit
is carried by B1\_07, B1\_06 and B2\_03, where the network loses 9.8--12.9 points. Power-MF was run
at our resampled rate with our matcher, so this is not the authors' own evaluation. The 10--60\Hz\ band is not our finding but Xu et al.'s \cite{xu2026}; our contribution
is only the measurement of what it is worth. That measurement (tier 1, $+$11.00) was made with a
gradient-boosted learner on 300\ms\ windows; on the TCN it falls to $+$2.44 with an
interval containing zero, so tier 1 is not established on the network and $+$11.00 must be read as
bounded to the 300\ms\ learner. That run is preliminary: one seed. B1 labels are indirect, so B1 F1 measures agreement with the authors'
pipeline; a trace of that provenance shows in timing: on the eight B1 women where all three models
exceed 95 F1 the mean absolute offset is 3.25\ms\ for the model trained on B1 against 6.50\ms\ for
models that never saw it, so B1 must not be used to argue timing accuracy. DPSS \cite{shokouhmand2023} numbers remain quoted, not re-run. The main model uses one
seed and the CinC figures one in-domain threshold per model. No subject lies in the range where an STV
threshold would change management, so we report a resolution failure, not a measured sensitivity. One topological configuration was tried, without tuning on test labels.

\section{Conclusion}
With protocol held fixed, single-channel fQRS detection on ADFECGDB reaches 99.21\% F1 on a
label-blind lead and 93.30\% zero-shot on pregnancy recordings. Re-run on all 22 women, the
four-channel Power-MF pipeline gains 12.12 F1 points over its own single-channel restriction; the
network recovers 10.85 of them, 89.5\%, from one lead, which leaves it not separable from four
channels ($-$1.27, CI $[-3.08;+0.27]$) and above one channel ($+$10.85, 22/22 women). The measurable
part of the remaining gain comes from 1.5\,s of temporal context; the 10--60\Hz\ front-end of
\cite{xu2026} is worth 11.00 points to a 300\ms\ learner but only 2.44,
preliminarily, to the network; the network family is a variable five women cannot resolve. Cross-recorder transfer, not gestational age, is the open
problem: over all 75 CinC 2013 set-a records the blind score is 79.40\% after training on 22 women,
and the lead rule degrades before the detector does. Derived short-term variability inherits
detection error one-sidedly and is not yet a clinical measurement.

\bibliographystyle{vancouver}
\scriptsize
\bibliography{refs}
\end{document}
